% Source: https://artofproblemsolving.com/community/c3h2351605_mathematics
% Stable identity: aops-topic-2351605-q01
Write $A=\overline{abcd}=1000a+100b+10c+d$. The new number is
\[1000(a+n)+100(b-n)+10(c+n)+(d-n)=A+909n.\]
So $A+909n=nA$, that is, $(n-1)A=909n$. Here $n\ne1$ (otherwise $909=0$), so $A=\dfrac{909n}{n-1}$. Since $n$ and $n-1$ have no common factor, $n-1$ must divide $909=3^2\cdot101$, so $n-1\in\{1,3,9,101,303,909\}$:
\begin{itemize}
\item $n=2$: $A=1818$;
\item $n=4$: $A=1212$, but its second digit $2$ cannot be decreased by $4$;
\item $n=10$: $A=1010$, but its first digit $1$ cannot be increased by $10$;
\item $n=102$, $304$, $910$: $A=918$, $912$, $910$, which have only three digits.
\end{itemize}
So $A=1818$ and $n=2$: indeed the digits become $3,6,3,6$, and $3636=2\cdot1818$. The sum of the digits of $A$ is $1+8+1+8=\boxed{18}$.
